\section{Aplicación de Git con Conventional Commits}

\subsection{Breve explicación de qué son los Conventional Commits y por qué se usan}
[Explica en tus propias palabras qué es la convención de commits (Conventional Commits) y las ventajas de utilizarla en el desarrollo]

\subsection{Guía de tipos principales}
[Proporciona una breve descripción de los tipos principales de commits que utilizarán en su proyecto. Ejemplo:]
\begin{itemize}
    \item \texttt{feat}: [Descripción de feat]
    \item \texttt{fix}: [Descripción de fix]
    \item \texttt{docs}: [Descripción de docs]
    \item \texttt{test}: [Descripción de test]
    \item \texttt{chore}: [Descripción de chore]
    \item \texttt{ci}: [Descripción de ci]
    \item \texttt{refactor}: [Descripción de refactor]
\end{itemize}

\subsection{Capturas del historial de commits}
[Inserta aquí las capturas de pantalla de su historial de commits (mínimo 9 registros). Luego, realiza un análisis detallado de una muestra de 6 commits.]

\begin{figure}[htbp]
    \centering
    % \includegraphics[width=0.8\textwidth]{figures/commits_history.png}
    \caption{Historial de Commits}
    \label{fig:commits}
\end{figure}

[Análisis detallado de los 6 commits seleccionados explicando el porqué de cada tipo y su impacto]

\subsection{Captura de Pull Requests}
[Inserta aquí la captura de los Pull Requests realizados siguiendo su metodología de branching (GitFlow o Trunk Based Development)]

\begin{figure}[htbp]
    \centering
    % \includegraphics[width=0.8\textwidth]{figures/pull_requests.png}
    \caption{Pull Requests siguiendo la metodología de ramificación}
    \label{fig:pull_requests}
\end{figure}
